\section{Canonical monomials and a new constraint}
\label{sec:main}

In this section we introduce the notion of ``canonical monomials''
of a given degree---%
simple %
monomials that appear in every affine-invariant
linear property containing monomials of a given degree.
We then give a constraint that rejects canonical monomials of
some special degrees, while accepting all monomials of lower degrees.

Later, in \expref{Section}{sec:final}, we show how to use this to build
a constraint whose orbit accepts all monomials of total degree at most $d$ while rejecting
all monomials of total degree $b_i(d)$, which suffices to get \expref{Theorem}{thm:main2}.

\begin{definition}[Canonical monomials]\label{defn:canonical-monomial}
Let $q=p^s$ for a prime $p$. The canonical monomial of (total) degree $d$ over $\F_q$ is the monomial $\prod_{i=1}^{\ell} x_i^{d_i} $
which satisfies $\sum_{i=1}^{\ell} d_i = d$,
$d_i = q-q/p$ for all $2 \leq i \leq \ell$,  $0 \leq d_1 \leq q-1$ and  $d_1 + q-q/p > q -1$.
\end{definition}

We note that~\cite{HSS11} used a different canonical monomial (cf.\ Definition 4.1.,~\cite{HSS11}) for the proof of their improved bounds on testing Reed-Muller codes.
Our different choice of canonical monomial is needed to construct single-orbit characterizations which improve on those given in~\cite{HSS11}
in terms of the number of queries.
The main property of the canonical monomial that we will use in
\expref{Section}{ssec:complete}
to prove \expref{Theorem}{thm:main2} is that
every affine-invariant linear family that contains any monomial
of total degree $d$ also contains the canonical monomial of degree $d$.
(See \expref{Lemma}{lem:canonical-monomial}.) 
This will imply in turn that if we can find constraints
that \emph{reject} this canonical monomial their orbit will reject every monomial
of total degree $d$.

\subsection{A new constraint on monomials of total degree
  \texorpdfstring{$< p(q - q/p)$}{< p(p - q/p)}}\label{sec:main-lemma}

The main technical novelty in our paper is a $k$-constraint $C$ 
that accepts all monomials of total degree strictly less than
$p(q-q/p)$ in $p$ variables but rejects the canonical monomial of degree $p(q-q/p)$
(note that the latter monomial also has $p$ variables) for
$k= (2^{p-1}+p-1)q^{p-1}$. We state the lemma below and devote the
rest of this section to proving this lemma.

\begin{lemma}[Main technical lemma]\label{lem:main-technical}
For every $q$ which is a power of a prime $p$ there exists a
$k$-constraint $C$ on $\set{\F_q^n \to \F_q}$ which accepts all monomials of total degree smaller than $p(q-q/p)$ in $p$
variables and rejects the canonical monomial $\prod_{i=1}^{p} x_i^{q-q/p}$
of degree $p(q-q/p)$  over $\F_q$, where $k = (2^{p-1}+p-1)q^{p-1}$.
\end{lemma}


It will be convenient for us to represent the constraint $C$ as a $p$-variate polynomial over $\F_q$. More precisely, suppose that
 $g(x)$ is a $p$-variate polynomial $g(x) \in \F_q [x_1,x_2,\ldots,x_p]$ that is non-zero on at most $k$ points in $\F_q^p$.
We associate with $g(x)$ the $k$-constraint $C=(\overline \alpha, \overline \lambda)$,
$\overline \alpha = (\alpha_1,\ldots,\alpha_k) \in (\F_q^p)^k$,  $\overline \lambda = (\lambda_1,\ldots,\lambda_k) \in \F_q^k$, where the vector $\overline \alpha$ consists
of all points  in $\F_q^p$ on which $g(x)$ is non-zero and $\lambda_j = g(\alpha_j)$ for all $ 1 \leq j \leq k$. Clearly, for every function
$f: \F_q^p \to \F_q$ it holds that
\begin{equation*}
\sum_{j=1}^k \lambda_j f(\alpha_j) = \sum_{\beta_1, \ldots, \beta_p \in \F_q} g(\beta_1, \ldots, \beta_p) \cdot f(\beta_1, \ldots, \beta_p)\,.
\end{equation*}

Thus for our purposes it suffices to find
%
%
a $p$-variate polynomial  $g(x) \in \F_{q} [x_1,x_2,\ldots,x_p]$
 with at most $k$ non-zero points in $\F_q^p$ such that $$ \sum_{\beta_1,\ldots,\beta_p \in \F_q} g(\beta_1,\ldots,\beta_p) \cdot M(\beta_1,\ldots,\beta_p) = 0$$ for
 every monomial in $p$ variables of total degree smaller than $p(q-q/p)$ and
 $$ \sum_{\beta_1,\ldots,\beta_p \in \F_q} g(\beta_1,\ldots,\beta_p) \cdot M(\beta_1,\ldots,\beta_p) \neq 0$$ when $M(x)$ is the canonical monomial of degree $p(q-q/p)$.

We start by describing a polynomial $P(x)$ that will satisfy the conditions we expect in $g$ above. The best way to describe this polynomial is via the notion of \emph{directional derivatives}. Let $f: \F_q \to \F_q$ be a function. Define the
derivative of $f$ in direction $y \in \F_q$ as $f_y(x) = f(x+y) - f(x)$.
Define the iterated derivatives as
\[f_{y_1,\ldots,y_\ell}(x) = (f_{y_1,\ldots,y_{\ell-1}})_{y_\ell}(x) = \sum_{I \subseteq [\ell]} (-1)^{|I|+1} f\bigg(x + \sum_{i \in I}y_i\bigg) \,.\]

Let $f(x)$ be the polynomial $f(x)=x_p^{q-1}$. Our polynomial $P(x)$ will be defined as follows.
\begin{equation}\label{eq:P(x)}
P(x) = \frac {f_{x_1,\ldots,x_{p-1}}(x_p)} {x_1 \cdots x_{p-1}} =
\frac {\sum_{I \subseteq [p-1]} (-1)^{|I|+1} (x_p + \sum_{i \in I} x_i)^{q-1}} {x_1 \cdots x_{p-1}}\,.
\end{equation}

To see that $P(x)$ is indeed a polynomial we need to show that $f_{x_1,\ldots,x_{p-1}}(x_p)$ is divisible by $x_1 \cdots x_{p-1}$. 
This  follows from the following lemma.

\begin{lemma}
Let $f$ be a univariate polynomial over $\F_q$. Then the polynomial $f_{y_1,\ldots,y_{\ell}}(x)$, as a polynomial in variables $x,y_1,\ldots,y_\ell$, is divisible by $y_1,\ldots,y_{\ell}$.
\end{lemma}

\begin{proof}
The proof is by induction on $\ell$. 
For $\ell =1$, if we write $f(x) = \sum_{d} c_d x^d$ then we have
%
\begin{align*}
f_{y_1}(x) & =  f(x+y_1) - f(x) = \sum_{d} c_d (x+y_1)^d - \sum_{d} c_d x^d \\
 & =   \sum_{d} c_d \sum_{i=0}^d \binom {d} {i} x^{d-i} y_1^{i} - \sum_{d} c_d x^d = \sum_{d >0} c_d \sum_{i=1}^d \binom {d} {i} x^{d-i} y_1^{i}\,.
\end{align*}
%
Thus all monomials in $f_{y_1}(x)$ contain the variable $y_1$ and hence $f_{y_1}(x)$ is divisible by $y_1$.

Next assume that $g(x):=f_{y_1,\ldots,y_{\ell -1}}(x)$ is divisible by $y_1,\ldots,y_{\ell-1}$. We shall show that $f_{y_1,\ldots,y_{\ell}}(x)$ is divisible by $y_1,\ldots,y_{\ell}$. Let $g(x)=\sum_{d} c'_d x^d$.
\begin{align*}
f_{y_1,\ldots,y_{\ell}}(x) & =  g(x+y_{\ell}) - g(x) = \sum_{d} c'_d (x+y_\ell)^d - \sum_{d} c'_d x^d \\
 & =   \sum_{d} c'_d \sum_{i=0}^d \binom {d} {i} x^{d-i} y_\ell^{i} - \sum_{d} c'_d x^d = \sum_{d >0} c'_d \sum_{i=1}^d \binom {d} {i} x^{d-i} y_\ell^{i}\,.
\end{align*}
So all monomials in $f_{y_1,\ldots,y_{\ell}}(x)$ contain the variable $y_\ell$ and hence $f_{y_1,\ldots,y_{\ell}}(x)$ is divisible by $y_\ell$. Furthermore, by the induction hypothesis we have that $g(x)$ is divisible by $y_1,\ldots,y_{\ell-1}$, so all monomials in both $g(x+y_{\ell})$ and $g(x)$ are divisible by $y_1,\ldots,y_{\ell-1}$. Consequently, all monomials in  $f_{y_1,\ldots,y_{\ell}}(x)$ are divisible by $y_1,\ldots,y_{\ell-1}$, so   $f_{y_1,\ldots,y_{\ell}}(x)$ is divisible by $y_1,\ldots,y_{\ell-1}$.
\end{proof}

In order to prove our main technical \expref{Lemma}{lem:main-technical} it suffices to show that the number of non-zero points of $P(x)$ in $\F_{q}^p$ is at most
$(2^{p-1}+p-1)q^{p-1}$, that it accepts all monomials in $p$ variables of total degree smaller than $p(q-q/p)$, and that it rejects the canonical monomial of degree $
p(q-q/p)$. We prove these three claims in Lemmas~\ref{lem:tech-lem-non-zeros}, \ref{lem:tech-lem-accept} and~\ref{lem:tech-lem-reject} below, respectively. Given these three
lemmas our main technical \expref{Lemma}{lem:main-technical} is immediate. We start with bounding the number of non-zeros of $P(x)$.

\begin{lemma}\label{lem:tech-lem-non-zeros}
The number of non-zero points of $P(x)$ in $\F_{q}^p$ is at most $(2^{p-1}+p-1)q^{p-1}$.
\end{lemma}

\begin{proof}
Let $\overline \beta= (\beta_1,\ldots,\beta_p) \in \F_q^p$. Suppose first that all first $p-1$ coordinates of $\overline \beta$ are non-zero, so that the denominator of $P(\overline \beta)$ is non-zero. Suppose furthermore that all terms of the form $(\beta_p + \sum_{i \in I} \beta_i)^{q-1}$ in the numerator of $P(\overline \beta)$ are non-zero. Then in this case we have that
\[P(\overline \beta) = \frac {\sum_{ I \subseteq [p-1]} (-1)^{|I|+1} \cdot 1 } {\beta_1 \cdots \beta_{p-1}} =  -\frac {\sum_{i=0}^{p-1} \binom {p-1} {i} (-1)^{i} } {\beta_1 \cdots \beta_{p-1}}  = 0 \,.\]

Thus we have that whenever $\beta_1,\ldots,\beta_{p-1}$ are all non-zero, $P(\overline \beta)$ can be non-zero only if at least one of the terms of the form $(\beta_p + \sum_{i \in I} \beta_i)^{q-1}$ in its numerator equals zero. Note that each such term has exactly $q^{p-1}$ assignments in $\F_q^p$ that make it zero. Since the number of terms in the numerator is $2^{p-1}$,
the total number of points in $\F_q^p$ that satisfy that at least one of the terms in the numerator equals zero is at most
$2^{p-1}\cdot  q^{p-1}$.

Concluding, we have that there are at most $2^{p-1} q^{p-1}$ vectors $\overline \beta \in \F_q^p$ which satisfy that $\beta_1,\beta_2,\ldots,\beta_{p-1}$ are all non-zero and in addition $P(\overline \beta) \neq 0$. Since there are at most
$(p-1)q^{p-1}$ elements $\overline \beta \in \F_q^p$ in which at least one of the first $p-1$ coordinates is zero, we conclude that the number of elements $\overline \beta \in \F_q^p$ such that $P(\overline \beta) \neq 0$ is at most \[2^{p-1} q^{p-1} + (p-1)q^{p-1} =
(2^{p-1}+p-1)q^{p-1}\,.\qedhere\] 
\end{proof}

We shall show later (in \expref{Lemma}{lem:tech-lem-non-zeros-lb}) that the bound on the number of non-zero points of $P(x)$ obtained in \expref{Lemma}{lem:tech-lem-non-zeros} is essentially tight.
Next we show that the constraint $C$ associated with  $P(x)$ accepts all monomials in $p$ variables of total degree smaller than $p(q-q/p)$.
For this we need the following well-known fact.

\begin{claim}\label{claim:power-sum}
Let $q$ be a prime-power, and let $i$ be an integer 
%
from
$\{0,1,\ldots,q-1\}$. Then
\[\sum_{\beta \in \F_q} \beta^i = \begin{cases}
  -1 & \quad i = q-1\,, \\
  0 & \quad \text{otherwise.}
\end{cases}
 \]
\end{claim}

\begin{proof}
Recall that the multiplicative group of $\F_q^*$ is cyclic and let $\gamma$ be a generator of this group. Then for all $i \in \{1,\ldots,q-2\}$ we have
$$\sum_{\beta \in \F_q} \beta^i = \sum_{j=1}^{q-1} (\gamma^i)^j = \gamma^i \cdot \frac {(\gamma^i)^{q-1} - 1} {\gamma^i -1} =  \gamma^i \cdot \frac {1- 1} {\gamma^i -1} =0\,, $$
whereas for $i = q-1$ we have that
\[\sum_{\beta \in \F_q} \beta^{q-1} = \sum_{\beta \in \F_q^*} 1 = q-1 = -1\,,\]
and for $i=0$ we have
\[\sum_{\beta \in \F_q} \beta^{0} = \sum_{\beta \in \F_q} 1 = q = 0\,.\qedhere\]
\end{proof}

\begin{lemma}\label{lem:tech-lem-accept}
Let $C$ be the constraint associated with $P(x)$. Then $C$ accepts all monomials in $p$ variables of total degree smaller than $p(q-q/p)$.
\end{lemma}

\begin{proof}
Let $m$ be a monomial in $p$ variables of total degree $d < p (q-q/p)$. We shall show that $$\sum_{\beta_1,\ldots, \beta_p \in \F_q} m(\beta_1,\beta_2,\ldots,\beta_p)  \cdot m'(\beta_1,\beta_2,\ldots, \beta_p)=0 $$ for every monomial $m'$ in $P(x)$. This will show in turn that $$\sum_{\beta_1,\ldots,\beta_p \in \F_q} m(\beta_1,\ldots,\beta_p) \cdot P(\beta_1,\ldots,\beta_p) =0$$ and hence $C$ accepts $m$.

Let $m'$ be a monomial in $P(x)$. First note that all monomials in the numerator of $P(x)$ have total degree exactly $q-1$ and hence all monomials
in $P(x)$ have total degree $q-1-(p-1)=q-p$. Thus $m'$ has total degree $q-p$, and $m \cdot m'$ is a monomial of total degree at most
$$q-p +d <
q-p + p(q-q/p) = p(q-1)$$ (after reducing individual degrees of variables mod $q$ if necessary).  Since $m \cdot m'$ is a monomial in $p$ variables, by pigeonhole principle there exists a variable $x_i$ in $m \cdot m'$ of degree smaller than $q-1$. Without loss of generality suppose that $x_1$ has degree $d' < q-1$ in $m \cdot m'$, and let $m \cdot m' = x_1^{d'} \cdot
\prod_{i=2}^{p} x_i^{d_i}$.

Thus we have
%
%
\begin{equation*}
 \sum_{\beta_1,\ldots,\beta_p \in \F_q} (m \cdot m') (\beta_1,\ldots,\beta_p)   = 
 \sum_{\beta_1,\ldots,\beta_p \in \F_q}  \beta_1^{d'} \cdot  \prod_{i=2}^{p} \beta_i^{d_i}  
 =  \bigg(\sum_{\beta_1 \in \F_q} \beta_1^{d'}\bigg) \prod_{i=2}^p \bigg( \sum_{\beta_i \in \F_q} \beta_i^{d_i}\bigg)
  =  0\,,
  \end{equation*}
%
%
where the last equality follows from \expref{Claim}{claim:power-sum} above and the fact that $d' <q-1$.
\end{proof}

We complete the proof of \expref{Lemma}{lem:main-technical} by showing that the constraint $C$ associated with $P(x)$ rejects the canonical
monomial of degree $p(q-q/p)$. In order to prove this, we shall use Kummer's Theorem which generalizes Lucas's Theorem (\expref{Theorem}{thm:luca}) and
gives a condition for when a multinomial coefficient
\[
\binom  {n}  {\gamma_1, \gamma_2,\ldots, \gamma_k} =
\frac {n!} {\gamma_1! \gamma_2! \cdots \gamma_k!} \not \equiv 0 \mod
p
\]
(it can be proved via repeated applications of Lucas's Theorem).


\begin{theorem}[Kummer's Theorem,~\cite{Kum}]
Let $n, \gamma_1,\gamma_2,\ldots,\gamma_k$ be integers such that $n = \gamma_1 + \gamma_2 + \cdots + \gamma_k$.
Then the multinomial coefficient 
\[
\binom {n}  {\gamma_1, \gamma_2,\ldots, \gamma_k} =
\frac {n!} {\gamma_1! \gamma_2 ! \cdots \gamma_k!} \not \equiv 0 \mod
p
\]
if and only if $\gamma_1, \gamma_2 , \ldots, \gamma_k$ sum
to $n$ in base-$p$ without carry.
\end{theorem}

\begin{lemma}\label{lem:tech-lem-reject}
Let $C$ be the constraint associated with $P(x)$. Then $C$ rejects the canonical monomial of degree $p(q-q/p)$ over $\F_q$.
\end{lemma}

\begin{proof}
The canonical monomial of degree $p(q-q/p)$ over $\F_q$ is the monomial $m = \prod_{i=1}^{p} x_i^{q-q/p}$. Let $m' = \prod_{i=1}^p x_i^{q/p-1}$.
We claim that
\begin{equation}\label{eq:eq1}
 \sum_{\beta_1,\ldots,\beta_p \in \F_q} (m \cdot m')(\beta_1,\ldots, \beta_p)  \neq 0\,,
 \end{equation}
 while for every other
monomial $m'' \neq m'$ in the support of
$P(x)$ it holds that
\begin{equation}\label{eq:eq2}
\sum_{\beta_1,\ldots,\beta_p \in \F_q}(m \cdot m'')(\beta_1,\ldots,\beta_p) =0\,.
\end{equation}
Thus in order to prove the lemma it will suffice to show that the monomial $m'$
is in the support of $P(x)$.
We start by showing \eqref{eq:eq1}. %
$$
\sum_{\beta_1,\ldots,\beta_p \in \F_q} (m \cdot m')(\beta_1,\ldots, \beta_p)   =  \sum_{\beta_1,\ldots,\beta_p \in \F_q}    \prod_{i=1}^{p} \beta_i^{q-1}
  = \prod_{i=1}^p \bigg( \sum_{\beta_i \in \F_q} \beta_i^{q-1} \bigg) =  (-1)^p  \,,
$$
where the last equality follows from \expref{Claim}{claim:power-sum} above.


Next we show that \eqref{eq:eq2} holds for every monomial $m'' \neq m'$ in the support of $P(x)$.
%
Let $m'' \neq m'$ be a monomial in the support of $P(x)$. Then
 $m''$ is a monomial of total degree $q-p$ and the fact that $m'' \neq m'$ implies that $m'' $ has a variable of degree smaller than
 $(q-p)/p = q/p - 1$. Without loss of generality suppose that the variable $x_1$ has degree smaller than  $q/p - 1$ in $m''$ and note that this implies
 that the variable $x_1$  has degree $d' < q-1$ in $m \cdot m''$. Let $m \cdot m'' = x_1^{d'} \cdot \prod_{i=2}^{p} x_i^{d_i}$. Then we have
 %
%
 %
%
%
%
\begin{equation*}
 \sum_{\beta_1,\ldots,\beta_p \in \F_q} (m \cdot m'')(\beta_1,\ldots, \beta_p)    =\sum_{\beta_1,\ldots,\beta_p \in \F_q}   \beta_1^{d'}  \prod_{i=2}^{p} \beta_i^{d_i} 
 =  \bigg(\sum_{\beta_1 \in \F_q} \beta_1^{d'} \bigg) \prod_{i=2}^p  \bigg(\sum_{\beta_i \in \F_q} \beta_i^{q-1} \bigg)= 0\,, 
\end{equation*}
where the last equality holds since $\sum_{\beta_1 \in \F_q} \beta_1^{d'}=0$ due to \expref{Claim}{claim:power-sum} above and the fact that $d' < q-1$.

%
 %
%
 It remains to show that
the monomial $m'$ is in the support of $P(x)$. This happens in turn if
and only if the monomial
\[
{\widetilde m} = x_1^{q/p-1} \prod_{i=2}^p x_i^{q/p}
\]
is in the numerator of $P(x)$.
Note that of all terms of the form $(x_p + \sum_{i \in I}x_i)^{q-1}$ in the numerator of $P(x)$, the monomial $\widetilde m$
can only belong to the support of $$\bigg(x_p + \sum_{i \in [p-1]} x_i\bigg)^{q-1} = (x_1 + x_2 + \cdots + x_p)^{q-1}\,.$$ Thus it suffices
to show that the monomial $\widetilde m$ belongs to the support of $(x_1 + x_2 + \cdots + x_p)^{q-1}$.

In order to show the above we resort to Kummer's Theorem. Expanding $(x_1+x_2+ \cdots + x_p)^{q-1}$ we have that the coefficient of the monomial $\widetilde m$ is
$\binom {q-1}  {\gamma_1, \ldots, \gamma_p}$ for $\gamma_1 = q/p -1$ and $\gamma_i = q/p$ for all $2 \leq i \leq p$. Noting that $\gamma_1, \ldots, \gamma_p$ sum to $q-1$ without carry in base-$p$, Kummer's Theorem implies that $\binom {q-1}  {\gamma_1, \ldots, \gamma_p}$ is non-zero mod $p$. This implies in turn that $\widetilde m$ is contained in the support of the polynomial $(x_1+x_2+ \cdots + x_p)^{q-1}$, thus completing the proof of the lemma.
 \end{proof}

 Given Lemmas~\ref{lem:tech-lem-non-zeros}, \ref{lem:tech-lem-accept} and~\ref{lem:tech-lem-reject} the proof of \expref{Lemma}{lem:main-technical} is immediate.

 \begin{proof}[Proof of \expref{Lemma}{lem:main-technical}]
Let $P(x)$ be the polynomial given in \eqref{eq:P(x)}, and let $C$ be the constraint on $\set{\F_q^p \to \F_q}$ associated with $P(x)$. {}From \expref{Lemma}{lem:tech-lem-non-zeros}
we have that the number of non-zero points of $P(x)$ in $\F_q^p$ is at most $(2^{p-1} + p-1)q^{p-1}$, and hence $C$ is a
($(2^{p-1}+p-1)q^{p-1}$)-constraint. \expref{Lemma}{lem:tech-lem-accept} implies that $C$ accepts all monomials of total degree smaller
than $p(q-q/p)$, while \expref{Lemma}{lem:tech-lem-reject} implies that $C$ rejects the canonical monomial of degree $p(q-q/p)$.
 \end{proof}

\subsection{Tightness of our analysis}

Next we show that the upper bound on the number of non-zero points of $P(x)$ in $\F_q^p$ given in \expref{Lemma}{lem:tech-lem-non-zeros}
is essentially tight.

\begin{lemma}\label{lem:tech-lem-non-zeros-lb}
The number of non-zero points of $P(x)$ in $\F_q^p$ is at least \[(2^{p-1}-p-1)q^{p-1} - 2^{p-1}(2^{p-1}-1)q^{p-2} \,.\]
\end{lemma}

\begin{proof}
We will show that the numerator of $P(x)$ is non-zero for at least $$2^{p-1}q^{p-1} -  2^{p-1}(2^{p-1}-1)q^{p-2}$$ elements in $\F_q^p$.
Since the denominator of $P(x)$ is zero for at most $(p-1)q^{p-1}$ elements in $ \F_q^p$ this will show that $P(x)$ is non-zero
for at least $$2^{p-1}q^{p-1} -  2^{p-1}(2^{p-1}-1)q^{p-2} - (p-1)q^{p-1}  $$ points in $\F_q^p$.

Let $ \overline \beta = (\beta_1,\ldots, \beta_p) $ be a random point in $\F_q^p$.
Let $E$ be the event that the numerator of $P(\overline \beta)$ is non-zero. Our goal will be to show that  $$\Pr[E] \geq  2^{p-1}/q - 2^{p-1}(2^{p-1}-1)/q^2\,.$$
For a subset $I \subseteq [p-1]$ let
$E_{I}$ be the event that $ (\beta_p + \sum_{i \in I} \beta_i )^{q-1} = 0$. Let $E'$ be the event that exactly one of the events $E_I$ holds. Clearly,
$E' \subseteq E$ and hence it suffices to show that $$Pr[E'] \geq  2^{p-1}/q - 2^{p-1}(2^{p-1}-1)/q^2\,.$$

 Note that $\Pr[E_I] = 1/q$ for all $I \subseteq [p-1]$ and
$\Pr[E_I \cap E_J] = 1/q^2$ for all $I \neq J$, $I,J \subseteq [p-1]$
since $\beta_p + \sum_{i \in I} \beta_i =0 $ and
$\beta_p + \sum_{j \in J} \beta_j = 0$ are linearly independent linear equations.

Thus we have that
%
\begin{align*}
\Pr[E'] &   \geq  \sum_{I \subseteq [p-1]} \bigg(\Pr[E_I] - \sum_{J \subseteq [p-1], J \neq I} \Pr[E_I \cap E_J] \bigg) \\
& = \sum_{I \subseteq [p-1]} \bigg(1/q -  \sum_{J \subseteq [p-1], J \neq I} 1/q^2 \bigg) \\[1ex]
& =   2^{p-1}(1/q - (2^{p-1} - 1)/q^2)  \\[1ex]
& =  2^{p-1}/q - 2^{p-1}(2^{p-1}-1)/q^2.
%
\end{align*}
\end{proof}



\section{Proof of \expref{Theorem}{thm:main2}}
\label{sec:final}

In this section we use \expref{Lemma}{lem:main-technical}
to prove
\expref{Theorem}{thm:main2}.
This part is done in several steps.
In \expref{Section}{ssec:arbit degree}
we extend the constraint from the previous section
to obtain constraints rejecting
canonical monomials of degree $d+1$, while accepting all monomials of total degree at most $d$ for an arbitrary integer $d$.
Next, in \expref{Section}{ssec:simult} we combine the various constraints from the previous
step to find one constraint which rejects all the canonical
monomials of degree $b_i(d)$ for every $0\leq i\leq s$, while accepting all
monomials of total degree at most $d$, for an arbitrary integer $d$. Finally, in \expref{Section}{ssec:complete}, we
prove \expref{Theorem}{thm:main2}. In this part, we use
some of the standard facts about affine-invariance to conclude that
the orbit of the constraint from the previous step must reject \emph{all monomials} (and not just the canonical monomials) of total degree $b_i(d)$ for $0\leq i\leq s$ while accepting all monomials of total degree at most $d$.
%
%

\subsection{Rejecting canonical monomials of arbitrary degree
  \texorpdfstring{$d +1$}{d + 1}}
\label{ssec:arbit degree}

We start by showing how the constraint guaranteed by
\expref{Lemma}{lem:main-technical} can be turned, via an operation on
 constraints that we call the \emph{product operation},
 into a constraint which rejects the canonical monomial of degree $d+1$ and accepts all monomials of total degree at most $d$ for an arbitrary integer $d$. This step is given in the following lemma.

\begin{lemma}\label{lem:arbitrary-degree}
For every  $q$ which is a power of a prime $p$, and 
%
for every pair $d$, $n$ of integers %
there exists a $k$-constraint $C$
on $\{\F_q^n \to \F_q\}$ which accepts all monomials
of total degree at most $d$
and rejects the canonical monomial
of degree $d+1$
over $\F_q$, where  \[k \leq q^2 \cdot
\left(2^{p-1}+p-1\right)^{(d+1)/(q(p-1))}
\cdot q^{(d+1)/q}\,.\]
\end{lemma}

For the proof of the above lemma first introduce the product operation. For a pair of vectors $\gamma=(\gamma_1,\ldots,\gamma_{n_1})$ and $\gamma'=(\gamma'_1,\ldots,\gamma'_{n_2})$ let $\gamma \circ \gamma'=
(\gamma_1,\ldots,\gamma_{n_1},\gamma'_1,\ldots,\gamma'_{n_2})$ denote their concatenation.


\begin{definition}[Product of constraints]
Let $$C_1=\bigg(\overline \alpha^{(1)}, \set{\overline \lambda^{(1)}_i}_{i=1}^{r_1}\bigg)$$ be a $k_1$-constraint
on $\set{\F_q^{n_1} \to \F_q}$,
and let  $$C_2=\bigg(\overline \alpha^{(2)}, \set{\overline \lambda^{(2)}_i}_{i=1}^{r_2}\bigg)$$ be a $k_2$-constraint on $\set{\F_q^{n_2} \to \F_q}$.
Their product  $C = C_1 \otimes C_2$ is the ($k_1 \cdot k_2$)-constraint  $$C=\bigg(\overline \alpha, \set{\overline \lambda_{(i_1,i_2)}}_{i_1=1,i_2=1}^{r_1,r_2}\bigg)$$
on $\set{\F_q^{n_1+n_2} \to \F_q}$, where
$\overline \alpha \in (\F_{q}^{n_1+n_2})^{k_1 \times k_2}$ and $\overline \lambda_{(i_1,i_2)} \in (\F_q)^{k_1 \times k_2}$ for all $1 \leq i_1 \leq r_1$, $1 \leq i_2 \leq r_2$,  and are defined as follows:
\[\alpha_{(j_1,j_2)} =  \alpha^{(1)}_{j_1} \circ  \alpha^{(2)}_{j_2}  , \; \lambda_{(i_1,i_2),(j_1,j_2)} = \lambda^{(1)}_{i_1,j_1}\cdot \lambda^{(2)}_{i_2,j_2} \]
for all $1 \leq j_1 \leq k_1$, $1 \leq j_2\leq k_2$, $1 \leq i_1 \leq r_1$, $1 \leq i_2 \leq r_2$.
\end{definition}

\begin{proposition}\label{pro:product}
Let $C_1$ be a $k_1$-constraint on $\set{\F_{q}^{n_1}\to \F_q}$ which accepts all monomials $x^{\overline d}$ with  $\overline d \in D_1$ and rejects all monomials $x^{\overline b}$ with $\overline b \in B_1$, and
let $C_2$ be a $k_2$-constraint on $\set{\F_{q}^{n_2}\to \F_q}$ which accepts all monomials $x^{\overline d}$ with  $\overline d \in D_2$ and rejects all monomials
$x^{\overline b}$ with  $\overline b \in B_2$.
Then $C_1 \otimes C_2$  is a $(k_1\cdot k_2)$-constraint on $\set{\F_q^{n_1+n_2} \to \F_q}$ which accepts
all monomials $x^{\overline d_1 \circ \overline e_2}$ and $x^{\overline e_1 \circ \overline d_2}$ where
$\overline d_1 \in D_1$, $\overline d_2 \in D_2$ and $\overline e_1 \in \{0,\ldots,q-1\}^{n_1}$, $\overline e_2 \in \{0,\ldots,q-1\}^{n_2}$ are arbitrary,
and rejects all monomials of the form
$x^{\overline b_1 \circ \overline b_2}$ where $\overline b_1 \in B_1$ and $\overline b_2 \in B_2$.
\end{proposition}

\begin{proof}

Let $\overline f = \overline f_1 \circ \overline f_2$ where $\overline f_1 \in \{0,\ldots,q-1\}^{n_1}$, $\overline f_2 \in \{0,\ldots,q-1\}^{n_2}$.
Then for every $1 \leq i_1 \leq r_1$, $1 \leq i_2 \leq r_2$ we have that

\begin{align*}
\sum_{j_1=1}^{k_1} \sum_{j_2=1}^{k_2} \lambda_{(i_1,i_2),(j_1,j_2)} (\alpha_{(j_1,j_2)})^{\overline f} & = 
\sum_{j_1=1}^{k_1} \sum_{j_2=1}^{k_2}  \lambda^{(1)}_{i_1,j_1} \cdot \lambda^{(2)}_{i_2,j_2} \cdot \big(\alpha^{(1)}_{j_1} \circ  \alpha^{(2)}_{j_2}\big)^{\overline f} \\
& = \sum_{j_1=1}^{k_1} \sum_{j_2=1}^{k_2}  \lambda^{(1)}_{i_1,j_1} \cdot \lambda^{(2)}_{i_2,j_2} \cdot \big(\alpha^{(1)}_{j_1}\big)^{\overline f_1}  \big(\alpha^{(2)}_{j_2}\big)^{\overline f_2} \\
%
& =  \bigg( \sum_{j_1=1}^{k_1}\lambda^{(1)}_{i_1,j_1}  \big(\alpha^{(1)}_{j_1} \big)^{\overline f_1}\bigg) \cdot  \bigg( \sum_{j_2=1}^{k_2}\lambda^{(2)}_{i_2,j_2}  \big(\alpha^{(2)}_{j_2}\big)^{\overline f_2}\bigg).
\end{align*}

Hence $C$ accepts all monomials of the form $x^{\overline d_1 \circ  \overline e_2}$ and $x^{\overline e_1 \circ \overline d_2}$ where
$\overline d_1 \in D_1$, $\overline d_2 \in D_2$ and $\overline e_1 \in \{0,\ldots,q-1\}^{n_1}$, $\overline e_2 \in \{0,\ldots,q-1\}^{n_2}$ are arbitrary and
 rejects all monomials of the form $x^{\overline b_1 \circ \overline b_2}$ where $\overline b_1 \in B_1$, $\overline b_2 \in B_2$.
\end{proof}

The product operation, applied to the constraint given by \expref{Lemma}{lem:main-technical}, suffices for proving \expref{Lemma}{lem:arbitrary-degree} when $p(q-q/p)$ divides $d+1$. However, since this is not always the case we need to consider also testing univariate monomials of degree at most $q-2$. The following lemma covers this case.

%
%

\begin{lemma}[Testing univariate monomials of degree at most $q-2$, (cf.~\cite{RS96})]\label{lem:constraint-uni}
Let $d$ be an integer %
from
$\{0,1,\ldots ,q-2\}$. Then there exists a ($d+2$)-constraint $C$ on $\{\F_q \to \F_q\}$ which accepts all monomials
$x^e$ for $e \leq d$ and rejects the monomial $x^{d+1}$.
\end{lemma}

\begin{proof}
Let $C=(\overline \alpha,\overline \lambda )$ be the $(d+2)$-constraint defined as follows.
Let $\alpha_1,\ldots,\alpha_{d+2} \in \F_q$ be distinct elements
(note that they do exist since $d+2 \leq q$).
Let ${\overline \lambda} = (\lambda_1,\ldots,\lambda_{d+2}) \in \F_q^{d+2}$
be a non-zero vector satisfying 
\[
\sum_{i=1}^{d+2} \lambda_i
\alpha_i^\ell = 0\quad\text{for all}\quad \ell \in \{0,\ldots,d\}\,.
\]
Note that such a vector ${\overline \lambda}$ exists since each of
the constraints above is a homogenous linear constraint on ${\overline \lambda}$
and there are only $d+1$ such constraints and $d+2$ variables.
We claim that
\[
\sum_{i=1}^{d+2} \lambda_i \alpha_i^{d+1} \ne 0\,,
\]
since
if it were then ${\overline \lambda}$ would be in the null space
of the Vandermonde matrix $[\alpha_i^j]_{i=1,j=0}^{d+2,d+1}$.
Thus $C$ rejects $x^{d+1}$ while accepting $x^{e}$
for every $e \leq d$.
\end{proof}

Given \expref{Proposition}{pro:product} and \expref{Lemma}{lem:constraint-uni} we are ready to prove \expref{Lemma}{lem:arbitrary-degree}.

\begin{proof}[Proof of \expref{Lemma}{lem:arbitrary-degree}]
Write $d+1$ as $d+1 = r+ \ell (q-q/p) $ where $0 \leq r \leq q-1$ and $r+q-q/p > q-1$. Write $\ell = \ell' p + r'$ where $0 \leq r' < p$.
Let $C_1$ be the ($r+1$)-constraint guaranteed by \expref{Lemma}{lem:constraint-uni} for the degree $r-1$, and let $C_2$
be the ($q-q/p+1$)-constraint guaranteed by the same lemma for the degree $q-q/p-1$. Let $C_3$ be
the $k'$-constraint guaranteed by \expref{Lemma}{lem:main-technical}.
Finally, let $C$ be the ($(r+1)\cdot (q-q/p+1)^{r'} \cdot (k')^{\ell'}$)-constraint which is the
product of $C_1$ with $r'$ copies of $C_2$ and $\ell'$ copies of $C_3$. That is,
$C =C_1 \otimes C_2^{\otimes r'} \otimes C_3^{\otimes \ell'}$.
We claim that $C$ accepts all monomials of total degree at most $d$ and rejects the canonical monomial of degree $d+1$
(If $C=\big(\overline \alpha, \bset{\overline \lambda_i }_i)$ is  a constraint on $\{\F_q^{n'} \to \F_q\}$ for $n' <n$ then we extend $C$ to be a constraint on $\{ \F_q^n \to \F_q \}$ by concatenating sufficient number of $1$'s
to each element $\alpha_j$ in the vector $\overline \alpha$.)

To see this suppose first that $m=x_1^{d_1} \cdots x_{n}^{d_n}$ is a monomial of total degree at most $d$.
In this case we have that either $d_1<r$ or $d_i <q-q/p$ for some $2
\leq i \leq r'+1$ or 
\[
\sum_{j=r' + 2 + p(i-1)}^{r' +1+pi} d_j < p(q-q/p)
\]
for some
$1 \leq i \leq \ell'$.
%
%
{}From \expref{Proposition}{pro:product} this implies that
the constraint $C$ accepts the monomial $m$. Suppose on the other  hand that $m$ is
the canonical monomial of degree $d+1$.  Then in this case we have that all
of the variables $x_2,\ldots, x_{\ell+1}$ are of degree $q-q/p$ and the variable $x_{1}$ is of degree $r$. Hence \expref{Proposition}{pro:product} implies that $C$ rejects the monomial $m$.


Finally note that the locality of $C$ is $k=(r+1)\cdot (q-q/p+1)^{r'}
\cdot (k')^{\ell'}$ where
\[
k' = (2^{p-1}+p-1)q^{p-1}\,,\quad \ell' \leq \frac{d+1}{(q-q/p)p}\,,\quad
r \leq q-1\,,\quad\text{and}\quad r' \leq p\,.
\]
Using the following series of simplifications, we can bound $k$ as claimed.
\begin{align*}
k
& =  (r+1)\cdot (q-q/p+1)^{r'} \cdot ((2^{p-1}+p-1)q^{p-1})^{\ell'} \\
& \leq  q \cdot q^{r'} \cdot ((2^{p-1}+p-1)q^{p-1})^{\ell'} \\
& =  q \cdot (2^{p-1}+p-1)^{\ell'} \cdot q^{r' + (p-1)\cdot \ell'} \\
& \leq  q \cdot (2^{p-1}+p-1)^{\ell'} \cdot q^{((p-1)/p) \cdot \ell + 1}\\
& \leq  q^2 \cdot (2^{p-1}+p-1)^{(d+1)/((q-q/p)p)} \cdot q^{((p-1)/p) \cdot ((d+1)/(q-q/p))}\\
& =  q^2 \cdot (2^{p-1}+p-1)^{(d+1)/(q(p-1))} \cdot q^{(d+1)/q}\,.\qedhere
\end{align*}
\end{proof}

\subsection{Rejecting all canonical monomials in the border simultaneously}
\label{ssec:simult}

Next we show for every integer $d$ the existence of a $k$-constraint which accepts all monomials of total degree at most $d$
and rejects all the canonical monomials of degree $b_i(d)$ for $0 \leq i \leq s$
simultaneously.
For proving this we shall use the union operation on constraints
defined as follows. For an integer $k$, let $0_k$ denote the all-zeros vector of length $k$.

\begin{definition}[Union of constraints]\label{def:union}
Let $$C_1=\bigg(\overline \alpha^{(1)}, \set{\overline \lambda^{(1)}_i}_{i=1}^{r_1}\bigg)$$ be a $k_1$-constraint on $\set{\F_q^n \to \F_q}$ and let
$$C_2=\bigg(\overline \alpha^{(2)}, \set{\overline \lambda^{(2)}_i}_{i=1}^{r_2}\bigg)$$ be a $k_2$-constraint on $\set{\F_q^n \to \F_q}$.
Their union $C = C_1 \cup C_2$ is the ($k_1+k_2$)-constraint  $$C=\bigg(\overline \alpha, \set{\overline \lambda'^{(1)}_i}_{i=1}^{r_1} \cup \set{\overline \lambda'^{(2)}_i}_{i=1}^{r_2}\bigg)$$  on $\set{\F_q^n \to \F_q}$ defined by
\begin{align*}
\overline \alpha &= \overline \alpha^{(1)} \circ \overline \alpha^{(2)}\,,\\
\overline \lambda'^{(1)}_i &= \overline \lambda^{(1)}_i \circ 0_{k_2} \; \text{for all} \; 1 \leq i \leq r_1\,,\\
\overline \lambda'^{(2)}_i &=  0_{k_1} \circ \overline
\lambda^{(2)}_i\; \text{for all} \; 1 \leq i \leq r_2\,.
\end{align*}
\end{definition}
 
\begin{proposition}\label{pro:union}
Let $C_1$ be a constraint which accepts all monomials with degrees in
 $D_1$ and rejects all monomials with degrees in $B_1$, and
let $C_2$ be a constraint which accepts all monomials with degrees in $D_2$ and rejects all monomials with degrees in $B_2$.
Then $C_1 \cup C_2$ accepts all monomials with degrees in $D_1 \cap D_2$ and rejects all monomials with degrees in $B_1 \cup B_2$.
\end{proposition}

\begin{proof}
For every degree $\overline e$,

\[\sum_{j=1}^{k_1+k_2} \lambda'^{(1)}_{i,j} \alpha_j^{\overline e} =  \sum_{j=1}^{k_1} \lambda^{(1)}_{i,j} \big(\alpha^{(1)}_j\big)^{\overline e} \; \text{for all $1 \leq i \leq r_1$}\,, \]

and similarly

\[\sum_{j=1}^{k_1+k_2} \lambda'^{(2)}_{i,j} \alpha_j^{\overline e} =  \sum_{j=1}^{k_2} \lambda^{(2)}_{i,j} \big(\alpha^{(2)}_j\big)^{\overline e} \; \text{for all $1 \leq i \leq r_2$}\,.\]

Thus the constraint $C$ accepts monomials of degree $\overline e$ if and only if both $C_1$ and $C_2$ accept the monomial of degree $\overline e$.
Hence $C$ accepts all monomials with degrees in $D_1 \cap D_2$ and rejects
all monomials with degrees in $B_1 \cup B_2$.
\end{proof}



Given the above proposition we can now build a constraint which
accepts all monomials of total degree at most $d$ while rejecting
all the canonical monomials of degree $b_i(d)$ for $0 \leq i \leq s$.

\begin{lemma}
\label{lem:canonical-monomial-single-orbit}
Let $q=p^s$ for a prime $p$ and let $n, d$ be arbitrary positive integers. Recall the definition of the integers $b_i(d)$ given in \expref{Definition}{defn:bi}.
Then there exists a $k$-constraint $C$ on $\{\F_q^n \to \F_q\}$
which accepts all monomials
of total degree at most $d$ and rejects all canonical
monomials of degree $b_i(d)$ for $0 \leq i \leq s$, where
\[k \leq 3q^4 \cdot \big(2^{p-1}+p-1\big)^{(d+1)/(q(p-1))} \cdot q^{(d+1)/q}\,.\]
\end{lemma}


\begin{proof}
For all $0 \leq i \leq s$ let $C_i$ be the $k_i$-constraint which accepts all monomials of total degree at most $b_i(d)-1$
and
rejects the canonical monomial of degree $b_i(d)$ over $\F_q$ as guaranteed by \expref{Lemma}{lem:arbitrary-degree}
with $$k_i \leq q^2  \cdot (2^{p-1}+p-1)^{(d+1+q)/(q(p-1))} \cdot q^{(d+1+q)/q} \leq 3q^3 \cdot (2^{p-1}+p-1)^{(d+1)/(q(p-1))} \cdot q^{(d+1)/q}\,.$$
Let $C = \bigcup_{i=0}^{s}C_i$. Then $C$ is a $k$-constraint for $$k = \sum_{i=0}^{s} k_i \leq 3q^4 \cdot (2^{p-1}+p-1)^{(d+1)/(q(p-1))} \cdot q^{(d+1)/q}\,.$$
\expref{Proposition}{pro:union} implies that $C$ accepts all monomials of total degree at most $d$ and rejects all canonical monomials
of degree $b_i(d)$ for $0 \leq i \leq s$, giving the claimed assertion.
\end{proof}





\subsection{Rejecting all monomials in the border}
\label{ssec:complete}

In order to complete the proof of \expref{Theorem}{thm:main2}, we show that for the constraint from \expref{Lemma}{lem:canonical-monomial-single-orbit}
which accepts all monomials of total degree at most $d$ while rejecting
all the canonical monomials of degree $b_i(d)$,  its orbit must accept all monomials of total degree at most $d$ while rejecting \emph{all monomials} of total degree $b_i(d)$.
%
and thus satisfies the conditions of \expref{Theorem}{thm:main2}.

We start by stating a simple (but useful) property of canonical monomials.

\begin{lemma}\label{lem:canonical-monomial}
Let $m$ be a monomial of total degree $d$, and let $\cF$ be a linear affine-invariant family containing $m$. Then $\cF$ contains the canonical monomial of degree $d$ over $\F_q$.
\end{lemma}


For the proof of the above lemma we shall need the following claim.

\begin{claim}\label{lem:canonical-monomial-bivariate}
Let $q = p^s$ for a prime $p$, and
let $m = x_1^{d_1} x_2^{d_2}$ be a monomial such that $0 \leq d_1 \leq q-1$, $0 \leq d_2 \leq q-1$ and $k \leq_p d_2$. Let $\cF$ be an affine-invariant linear family containing $m$.
 Then
the monomial $m' = x_1^{d_1+ k} x_2^{d_2 - k}$ is contained in $\cF$.
\end{claim}

\begin{proof}
Let $T$ be the affine-transformation $T(x_1, x_2) = (x_1, x_1+x_2)$.
Then
\[m \circ T = x_1^{d_1} (x_1+x_2)^{d_2} = x_1^{d_1} \bigg( \sum_{i=0}^{d_2} \binom {d_2} { i} x_1^i x_2^{d_2 - i} \bigg) =
\sum_{i=0}^{d_2} \binom {d_2} {i} x_1^{d_1+i} x_2^{d_2 - i} \,.\]
{}From Lucas's Theorem (\expref{Theorem}{thm:luca}) we have that $\binom {d_2} {i} \not \equiv 0 \mod p$ if and only if $i \leq_p d_2$.
So all monomials of the form $x_1^{d_1+i} x_2^{d_2 - i}$ such that $i \leq_p d_2$ are contained in the support of
$m \circ T$. 
%
Since $\cF$ is affine-invariant we have that $m \circ T$ is contained in $\cF$ and hence the
%
Monomial Extraction Lemma %
(\expref{Lemma}{lem:monomial-extraction}) implies that $m'$ is contained in $\cF$.
\end{proof}



Next we prove \expref{Lemma}{lem:canonical-monomial} based on \expref{Claim}{lem:canonical-monomial-bivariate}.

\begin{proof}[Proof of \expref{Lemma}{lem:canonical-monomial}]
Write $d = \ell(q-q/p) + r'$ where $r' <q-q/p$. Let $m = \prod_{i=1}^{n} x_i^{d_i}$.
We start by showing that $\cF$ contains a monomial $\prod_{i=1}^{n} x_i^{d'_i}$ which satisfies $d'_i \geq q-q/p$ for every $2 \leq i \leq \ell+1$.
 %

We shall apply \expref{Claim}{lem:canonical-monomial-bivariate} iteratively. For a monomial $m' = \prod_{i=1}^n x_i ^{e_i}$
let  $$c(m') = \sum_{i=2}^{\ell+1} \max\{0,(q-q/p)-e_i\}\,.$$ Clearly, $c(m') =0$ if and only if $e_i \geq q-q/p$ for all $2 \leq i \leq \ell+1$.
If $m$ satisfies that $d_i \geq q-q/p$ for all $2 \leq i \leq \ell+1$
 then we are done, hence assume that there exists $2 \leq i \leq \ell+1$ such that $d_i < q-q/p$.
If there exists a degree $j \in \set{1}\cup \set{\ell+2,\ldots,n}$ such that $d_j >0$ let $p^k$ be such that $p^k \leq_p d_j$.
\expref{Claim}{lem:canonical-monomial-bivariate} implies that $m_1 := m\cdot x_i^{p^k} x_j^{-p^k}$ is contained in $\cF$.  Otherwise,
since $ \sum_{i=1}^{n} d_i  = d$ there exists $j \in \set{2,\ldots,\ell+1}$ such that $d_j > q-q/p$. Note that the base-$p$ representation of $q-q/p$
is $(q-q/p) =  (p-1)\cdot p^{s-1}$ and hence the fact that $d_j > q-q/p$ implies the existence of $p^k \leq_p d_j$ such that
$q-q/p + p^k \leq_p d_j$. \expref{Claim}{lem:canonical-monomial-bivariate} implies that $m_1 := m\cdot x_i^{p^k} x_j^{-p^k}$  is contained in $\cF$ in this case as well.
Note that in both cases we have that $c(m_1) < c(m)$. Also, since $d_i < q-q/p$ we have that $d_i + p^k < q-q/p + q/p <q$. Hence all variables in $m_1$ have degree at most $q-1$ (we mention this fact since this will allow us to apply \expref{Claim}{lem:canonical-monomial-bivariate} iteratively).

If all variables $x_2,\ldots,x_{\ell+1}$ in $m_1$ have degree at least $q-q/p$
then we are done. Otherwise repeat the same process as previously to obtain a monomial $m_2 \in \cF$
such that the degrees of all variables in $m_2$ are at most $q-1$ and  $c(m_2) < c(m_1)$.
Continuing this way we have that
at the $i$-th step either all variables $x_2,\ldots,x_{\ell+1}$ in $m_{i-1}$ have degree at least $q-q/p$ and hence we are done
 or that we obtain a monomial $m_i \in \cF$ such that the degrees of all variables
in $m_i$ are at most $q-1$ and $c(m_i) < c(m_{i-1})$. Since the function $c(m_i)$ strictly declines at each step the process
must terminate eventually, and when it terminates we have that  $m_{i}$ satisfies that all variables $x_2, \ldots,x_{\ell+1}$ in it have degree at least $q-q/p$.

We have just shown that $\cF$ contains a monomial $m' = \prod_{i=1}^n x_i^{d'_i}$ where $d'_i \geq q-q/p$ for all $2 \leq i \leq \ell+1$.
Next we claim that the monomial ${\widetilde m} = \prod_{i=1}^n x_i^{{\widetilde d}_i}$ which satisfies
${\widetilde d}_1 = r'$, ${\widetilde d}_i=q-q/p$ for all $2 \leq i \leq \ell+1$ and ${\widetilde d}_i =0$ for all $\ell+2 \leq i \leq n$ is contained in $\cF$ (note that $\widetilde m$ is the canonical monomial of degree $d$ if and only if $r' + q-q/p > q-1$). To see this note that if $m' = \widetilde m$ then we are done.
Otherwise we must have that $d'_1 < r'$. As was the case previously this implies the existence of either $\ell+2 \leq j \leq n$ and an integer $k$ such that $p^k \leq_p d'_j$
or $2 \leq j \leq \ell+1$ such that $q-q/p + p^k \leq_p d'_j$. \expref{Claim}{lem:canonical-monomial-bivariate} implies that the monomial $m' \cdot x_1^{p^k} x_j^{-p^k}$ is contained
in $\cF$. Note that in  the latter monomial the degree of the variable $x_1$ increased, but the degree of all variables $x_2,\ldots,x_{\ell+1}$ remained at least $q-q/p$. Continuing this way we conclude that
the monomial $\widetilde m $ is contained in $\cF$.

Finally, note that if $r' +q-q/p >q-1$ then $\widetilde m $ is also the canonical monomial
of degree $d$ and hence we are done. Otherwise we have that $q-q/p \leq_p {\widetilde d}_{\ell+1}$ and hence \expref{Claim}{lem:canonical-monomial-bivariate} implies that the monomial
$ {\widetilde m} \cdot x_1^{q-q/p} \cdot x_{\ell+1}^{-(q-q/p)}$ is contained in $\cF$. The proof is completed by noting that in this case the latter monomial is the canonical monomial of degree
$d$.
\end{proof}

We are now ready for the proof of \expref{Theorem}{thm:main2}.

\begin{proof}[Proof of \expref{Theorem}{thm:main2}]
{}From \expref{Lemma}{lem:canonical-monomial-single-orbit}
we have a constraint $C$ of arity
\[k \leq 3q^4 \cdot \big(2^{p-1}+p-1\big)^{(d+1)/(q(p-1))} \cdot q^{(d+1)/q}\] that
accepts every monomial of total degree at most
$d$ and rejects every canonical monomial of degree $b_i(d)$ for $0 \leq i \leq s$.

%
%
The constraint $C$ accepts all monomials
of total degree at most $d$ and hence $C$ accepts all functions in $\mrm[n,d,q]$.
Since $\mrm[n,d,q]$ is affine-invariant this implies in turn that the orbit of $C$ accepts all functions in $\mrm[n,d,q]$ as well, and in particular all monomials of total degree at most $d$.

It just remains to show that the orbit of $C$ rejects every monomial of total degree $b_i(d)$ for $0 \leq i\leq s$.
Assume for contradiction that the orbit of $C$ accepts a monomial $m$ of total degree $b_i(d)$ for some
$0 \leq i \leq s$.
Let $\calf'$ be the set of functions accepted by the orbit of
$C$, \ie,
$$\calf' = \{f : \F_q^n \to \F_q \;|\;
\mbox{$T \circ C$~accepts~$f$~for~every~affine-transformation~$T\}$}\,.$$
We have that $\calf'$ is linear and affine-invariant, and contains $m$,
and so by \expref{Lemma}{lem:canonical-monomial} it also contains
the canonical monomial of degree $b_i(d)$. So the orbit of $C$ accepts the canonical monomial of degree $b_i(d)$ contradicting the
hypothesis about $C$.

%
%
\end{proof}
